\section{Keep the continuum: response and numerical complex scaling}
\label{sec:continuum-response}
% BEGIN CURATED SUBJECT INDEX
\index{complex scaling}
% END CURATED SUBJECT INDEX

\subsection{Continuum level density}
\label{subsec:cld}
% BEGIN CURATED SUBJECT INDEX
\index{spectrum!point}
\index{resolvent}
\index{trace-class operator}
\index{cross section}
\index{continuum level density}
% END CURATED SUBJECT INDEX

The density of states of an unbounded system is infinite and depends on the
normalization volume.  A finite observable is obtained by subtracting a
reference Hamiltonian $H_0$ with the same asymptotic channels.  Define the
continuum level density
\begin{equation}
  \Delta\rho(E)
  =-\frac{1}{\pi}\im\Tr\left[
  \frac{1}{E+\rmi0-H}
  -\frac{1}{E+\rmi0-H_0}
  \right] .
  \label{eq:cld-physical}
\end{equation}
When the resolvent difference is trace class, analytic deformation gives
\begin{equation}
  \Delta\rho(E)
  =-\frac{1}{\pi}\im\Tr\left[
  \frac{1}{E-H_\theta}
  -\frac{1}{E-H_{0,\theta}}
  \right] .
  \label{eq:cld-csm}
\end{equation}
The right-hand side can be evaluated on the real energy axis because the
deformed poles and continuum lie below it.  In an exact calculation it is
independent of $\theta$ within the common analyticity domain
\cite{Myo:2014ypa,Morikawa:2025vvs,Ogawa:2026dSCSM}.

If a finite basis yields complex eigenvalues $E_n^{\theta}$ and reference
eigenvalues $E_{n,0}^{\theta}$, a useful approximation is
\begin{align}
  \Delta\rho_N(E)
  &=-\frac{1}{\pi}\im
  \left[
  \sum_{n=1}^{N}\frac{1}{E-E_n^{\theta}}
  -\sum_{n=1}^{N}\frac{1}{E-E_{n,0}^{\theta}}
  \right] .
  \label{eq:finite-cld}
\end{align}
The subtraction must use compatible bases, contours, and asymptotic
Hamiltonians.  Otherwise large discretization artifacts survive and can be
mistaken for response.

For one elastic channel, the spectral-shift relation is
\begin{equation}
  \Delta\rho(E)=\frac{1}{\pi}
  \frac{\rmd\delta(E)}{\rmd E} ,
  \label{eq:cld-phase}
\end{equation}
up to the convention for the reference phase and bound-state jumps.  More
generally, if $\det S(E)=\rme^{2\rmi\vsub{\delta}{tot}(E)}$, then
\begin{equation}
  \Delta\rho(E)
  =\frac{1}{2\pi\rmi}
  \frac{\rmd}{\rmd E}\log\det S(E)
  =\frac{1}{\pi}
  \frac{\rmd\vsub{\delta}{tot}(E)}{\rmd E} .
  \label{eq:birman-krein-differential}
\end{equation}
Thus the CLD retains both pole and nonresonant phase information.  It is not,
by itself, the transmission probability: channel-resolved amplitudes are
needed to reconstruct greybody factors or reaction cross sections.

\subsection{Basis expansion and generalized eigenvalue problem}
\label{subsec:csm-basis}
% BEGIN CURATED SUBJECT INDEX
\index{spectrum!point}
% END CURATED SUBJECT INDEX

Choose square-integrable basis functions $\{\phi_n\}_{n=1}^{N}$.  Expanding a
right eigenvector as $\Psi^{\mathrm R}=\sum_n c_n\phi_n$ gives
\begin{equation}
  \sum_{n=1}^{N}H_{mn}^{\theta}c_n
  =E\sum_{n=1}^{N}N_{mn}c_n ,
  \label{eq:generalized-eigenproblem}
\end{equation}
where
\begin{align}
  H_{mn}^{\theta}&=\cprod{\phi_m}{H_\theta\phi_n},
  &N_{mn}&=\cprod{\phi_m}{\phi_n} .
  \label{eq:csm-matrices}
\end{align}
For an orthonormal basis, $N$ is the identity.  Gaussian bases are convenient
because kinetic and many potential matrix elements are analytic.  Even and
odd sectors on the full line can be represented by
\begin{align}
  \phi_n^{(+)}(x)&=\rme^{-\alpha_nx^2},
  &\phi_n^{(-)}(x)&=x\rme^{-\alpha_nx^2},
  \label{eq:real-gaussian-basis}
\end{align}
where every width satisfies $\re\alpha_n>0$.  Complex-range Gaussians replace
$\alpha_n$ by conjugate pairs $\alpha_n\pm\rmi\beta_n$ and resolve oscillatory
tails more efficiently
\cite{Hiyama:2003cu,Myo:2014ypa,Ogawa:2026veu}.

The generalized eigenproblem is ill conditioned when basis functions become
nearly linearly dependent.  A robust implementation first diagonalizes or
factorizes the overlap matrix, removes directions below a documented
tolerance, and only then solves the reduced problem.  Since the matrix is
non-normal, residuals must be checked for both right and left eigenvectors:
\begin{align}
  \epsilon_n^{\mathrm R}
  &=\frac{\lVert H^\theta\bm c_n
  -E_nN\bm c_n\rVert_2}
  {\lVert H^\theta\rVert_2\lVert\bm c_n\rVert_2
  +|E_n|\lVert N\rVert_2\lVert\bm c_n\rVert_2},
  \notag\\
  \epsilon_n^{\mathrm L}
  &=\frac{\lVert (\bm d_n)^{\dagger}H^\theta
  -E_n(\bm d_n)^{\dagger}N\rVert_2}
  {\lVert H^\theta\rVert_2\lVert\bm d_n\rVert_2
  +|E_n|\lVert N\rVert_2\lVert\bm d_n\rVert_2} .
  \label{eq:left-right-residuals}
\end{align}
Here $\bm c_n$ and $\bm d_n$ are coefficient vectors, and
$\lVert\cdot\rVert_2$ is the Euclidean or induced matrix norm.

\subsection{Pole-identification protocol}
\label{subsec:pole-protocol}
% BEGIN CURATED SUBJECT INDEX
\index{spectrum!point}
\index{pseudospectrum}
\index{pseudostate}
% END CURATED SUBJECT INDEX

A reproducible complex-scaling calculation should separate theorem-level
invariance from finite-basis evidence.  We recommend the following tests.

First, track eigenvalue trajectories under a range of scaling angles.  A
continuum pseudostate follows the ray $\arg(E-E_a)\simeq-2\theta$, whereas a
resonance develops a stationary segment after exposure.  Second, repeat the
tracking while increasing $N$ and changing the distribution of basis widths.
Third, compare at least two basis families when possible.  Fourth, monitor the
condition number of $N$, left--right overlap, and residuals.  Fifth, inspect
the smallest singular value of $zN-H^\theta$ around a candidate pole.  A large
pseudospectral cloud warns that small discretization errors can move the
eigenvalue substantially even when its trajectory looks stationary.

For black-hole problems, a sixth test is available: compare low-lying
frequencies against a method with different analytic systematics, such as
Leaver's continued fraction
\cite{Leaver:1985ax,Ogawa:2026veu}.  Agreement of a few digits is not the end
of validation; the scaling-angle plateau and basis stability identify why the
number belongs to the intended pole rather than to a coincidentally accurate
pseudostate.

\subsection{Complex absorbing potentials and stabilization}
\label{subsec:cap-stabilization}
% BEGIN CURATED SUBJECT INDEX
\index{spectrum!point}
\index{complex scaling}
\index{complex absorbing potential}
\index{stabilization method}
% END CURATED SUBJECT INDEX

An alternative is to add a complex absorbing potential,
\begin{equation}
  H(\eta)=H-\rmi\eta W,
  \qquad \eta>0,
  \qquad W\succeq0,
  \label{eq:cap}
\end{equation}
with $W$ supported outside the interaction region.  Resonances are inferred
from stationary eigenvalue trajectories as $\eta$ varies.  The logic is close
to complex scaling, but a finite absorber changes the operator rather than
implementing an exact analytic deformation.  Results should be extrapolated
or optimized in $\eta$ and absorber shape.  Stabilization methods similarly
identify resonances by weak dependence on box size or nonlinear basis scales.
These methods are valuable cross-checks but do not replace a definition of
the continued pole.
